\documentclass[11pt]{article}

\usepackage[utf8]{inputenc}
\usepackage{amsmath}
\usepackage{amssymb}

\title{External Effect Anomalies in AI Agent Workflows}
\author{}
\date{}

\begin{document}

\maketitle

\section{Introduction}

AI agents increasingly take actions that change the outside world: sending
emails, charging cards, placing orders, calling external services. Such
effects often cannot be undone. Take an agent booking a trip on a 1500 EUR
budget. To save time it books the flight (900) and the hotel (700) in
parallel. Both calls succeed, but together they break the budget. The hotel
can be canceled; the flight was non-refundable. Every tool call worked, yet
the workflow left a state in the world that cannot be fully repaired.

Agents make this common because they run steps in parallel (supervisor-worker
decomposition), plan speculatively (tree search, best-of-$k$), execute likely
actions early to hide latency, and recover by retrying after failures. Each
is useful, but each can also produce a new kind of external-effect failure.
We argue these failures form one catalog, with three sources.
\emph{Uncertainty} failures: when the system cannot confirm the outcome of
an effect --- because of a crash, a timeout, or a lost acknowledgment --- and
must act anyway, the same effect may happen
twice, fail to happen while the workflow proceeds as if it had, or be
compensated blindly. \emph{Workflow} failures: an effect may outlive an
aborted workflow, leak before the workflow decides, or be depended on by a
committed branch though it came from an abandoned one. \emph{Interaction}
failures: an effect meets something outside its own execution --- a
concurrent execution writing the same resource, or an outside actor that
observes and reacts to it before it can be undone. Agents show all three,
because an agent is a long-running workflow that acts on tools whose
outcomes it cannot always confirm and shares an open world with others.

We define correctness through a catalog of anomalies, each a forbidden
pattern in the effect history. Three features shape the setting and
separate it from a classical transaction: there is no rollback (an external
effect can only be compensated, and compensation is optional and may fail);
effects may be released before the workflow decides; and the surrounding world
is open and may react.

\section{Two Levels of Agent Operations}

We distinguish two abstraction levels L0 and L1. This distinction is inspired by multilevel transaction management, where
high-level transactional guarantees are derived from the properties of
lower-level operations~\cite{weikum1991principles}. Agent workflows
exhibit a similar separation, but with a different boundary. In classical
systems, lower-level operations are typically under the control of the
database system, and their semantic properties can be derived from the
implementation. In agent systems, L0 operations are often black-box tools
behind external APIs. Their properties---such as idempotence,
invertibility, externalization, and determinism---must therefore be
declared, inferred, or validated empirically. These L0 properties then
constrain which effect-safety guarantees are achievable for the composed L1
workflow.

\subsection{L0: Tool Operations}

L0 consists of individual effectful operations, e.g.:

\begin{itemize}
    \item \texttt{charge()}
    \item \texttt{send\_email()}
    \item \texttt{book\_flight()}
    \item \texttt{publish\_post()}
\end{itemize}

Each operation possesses properties:

\begin{itemize}
\item idempotence,
\item invertibility,
\item externalization,
\item determinism,
\item commutativity.
\end{itemize}

These properties are independent: \texttt{increment} commutes but is not
idempotent, while \texttt{delete} is idempotent but does not commute. They
therefore form a multidimensional property space rather than a single scale
(e.g.\ a reversible/irreversible axis). Commutativity is moreover \emph{relational}
(a property of operation pairs) and often \emph{conditional} on state
(escrow-style~\cite{oneil1986escrow}); determinism, absent from classical
concurrency control, is an additional axis that black-box LLM-backed tools
routinely violate.

\subsection{L1: Workflow Operations}

L1 consists of agent tasks, e.g.:

\begin{itemize}
    \item \texttt{HireCandidate}
    \item \texttt{BookTrip}
    \item \texttt{NegotiateContract}
    \item \texttt{PurchaseEquipment}
\end{itemize}

Workflows turn individual tool calls into a logical unit of work. Tool properties constrain the guarantees achievable by workflows, but
effect safety itself is defined at L1.

\section{Effect Histories}

We model execution as histories consisting of:

\begin{itemize}
\item $effect(o)$: externalization of operation $o$ --- the moment its
effect becomes visible outside the runtime and is no longer under full
transactional control. When it matters, we write $effect(o \text{ on } r)$
to name the external resource $r$ the effect acts on,
\item $cmp(o)$: compensation,
\item $dep(o_2 \leftarrow o_1)$: $o_2$ causally depends on the effect of
$o_1$ (its decision was made observing $effect(o_1)$),
\item $commit(w)$,
\item $abort(w)$.
\end{itemize}

Both $effect(o)$ and $cmp(o)$ carry an outcome status: \emph{confirmed},
\emph{failed}, or \emph{unknown} (no acknowledgment was received). We
additionally define:

\begin{itemize}
\item $resolved(w)$: a $commit(w)$ or $abort(w)$ event has occurred, i.e.\
the workflow's outcome is decided and will not change.
\item $survives(o)$: $effect(o)$ may have occurred (\emph{confirmed} or
\emph{unknown}) and there is no later successful $cmp(o)$ --- the effect
remains in the external world.
\end{itemize}

Unlike classical read-write histories~\cite{adya1999},
effect histories describe what the workflow externalized to the
outside world. The interesting anomalies live in the window between
$effect$ and $resolved$/$cmp$ --- a window absent from read-write models.

\section{Effect Anomalies}

Each anomaly is a \emph{consequence of a missing property}. We state each in
one sentence and then give a formal trace as an illustrative example; the
trace is one way the anomaly can arise, not its definition. We tag each
anomaly by its \emph{source}: \emph{uncertainty} (A1--A3), \emph{workflow}
(A4--A6), or \emph{interaction} (A7--A8). The uncertainty anomalies arise when
the system cannot confirm the outcome of an effect and must act anyway; the
unconfirmed outcome can come from a crash, a timeout, or a lost or absent
acknowledgment. The workflow anomalies arise from the structure of a
single execution --- survival under abort, dependency, and timing relative to
its resolution. The interaction anomalies arise when an effect meets
something outside its own execution: another execution writing the same
resource (A7), or an external actor that observes and reacts to it (A8).

\paragraph{A1: Duplicated Effect (source: uncertainty)}
\emph{A single logical effect is externalized more than once.} This arises
when the outcome of a call is not confirmed, so the system re-issues it, and
both attempts take effect. For example:
\begin{verbatim}
effect(pay_invoice) = unknown      no confirmation (timeout)
reissue(pay_invoice)               assume first didn't happen
effect(pay_invoice) = confirmed    but the first succeeded too
=> invoice paid twice
\end{verbatim}
Non-determinism makes this worse: re-issuing a non-deterministic operation
yields not a duplicate but a \emph{different} effect. Idempotence neutralizes
repetition; determinism does not by itself make re-issue safe, but it makes
replay matching meaningful, since a re-issued call then denotes the same
effect.

\paragraph{A2: Missing Committed Effect (source: uncertainty)}
\emph{A workflow commits relying on an effect that did not successfully
externalize --- either because an unknown outcome was resolved optimistically
or because a failed or absent required effect was ignored.} The canonical
case is uncertainty resolved the wrong way: the system assumes an unconfirmed
effect succeeded and commits, though it had actually failed. For example:
\begin{verbatim}
effect(pay_invoice) = unknown    outcome not confirmed
commit(Order)                    assumes success (it failed)
=> required effect absent, yet the workflow committed
\end{verbatim}
The anomaly is narrow: it is \emph{not} that the agent omitted a step
(a planning failure), but that a required effect is absent from the
committed history --- whether its outcome was unknown and guessed wrong, or
known to have failed and committed over anyway.

\paragraph{A3: Orphaned Compensation (source: uncertainty)}
\emph{A compensation is issued for an effect whose outcome is unknown.} This
arises when an unconfirmed outcome is instead resolved pessimistically: the
system compensates without knowing whether the original effect happened. For
example:
\begin{verbatim}
effect(pay_invoice) = unknown      outcome not confirmed
cmp(pay_invoice)    = refund       refund issued blindly
=> if payment never happened, the refund is spurious
\end{verbatim}
The anomaly is the \emph{blind action} at the moment of $cmp$, not the
eventual outcome; the remedy is to resolve $unknown$ before compensating.
A1, A2, and A3 are a family: all three stem from an outcome the system
cannot confirm, and differ only in the reaction. Faced with an unconfirmed
outcome, the system may re-issue (and duplicate, A1), assume success and
commit (and miss a required effect, A2), or compensate blindly (A3).
Resolving the outcome before acting removes all three.

\paragraph{A4: Irreversible Residue (source: workflow)}
\emph{An aborted workflow leaves behind an effect that cannot be retracted.}
For example, a non-refundable ticket is purchased, then the trip is aborted
and the purchase cannot be undone:
\begin{verbatim}
effect(book_flight) = confirmed    non-refundable
abort(BookTrip)                    budget violated
cmp(book_flight)    = failed
=> survives(book_flight): cannot retract; money lost
\end{verbatim}

\paragraph{A5: Premature Externalization (source: workflow)}
\emph{An effect that may not survive is externalized before its workflow is
resolved.} For example, an offer letter is sent before legal approval
completes, and approval is later denied:
\begin{verbatim}
effect(send_offer) = confirmed     letter externalized
NOT resolved(Hire)                 not yet committed/aborted
abort(Hire)                        approval later denied
=> effect(send_offer) preceded resolved(Hire)
\end{verbatim}
Premature externalization separates compensation-safe execution, which may
externalize an effect early and compensate later, from speculation-safe
execution, which stages or gates effects that may not survive until the
workflow's resolution is known.

\paragraph{A6: Contaminated Speculation (source: workflow)}
\emph{A committed effect causally depends on an effect from a branch that
does not survive.} For example, a committed branch externalizes an action
based on a temporary reservation made by another branch that is later
canceled:
\begin{verbatim}
effect(reserve_room R) = confirmed [b1]  room R held
dep(send_invites <- reserve_room R) [b2] b2 needs R
effect(send_invites) = confirmed [b2]    invites naming R sent
commit(b2)                               b2 wins, committed
abort(b1); cmp(reserve_room R)           b1 loses; R released
=> b2 already announced R, which no longer survives
\end{verbatim}
The contamination is causal, not about reversibility: $b2$ trusted an
effect of an unresolved branch. This is the effect-world analogue of a
\emph{dirty read} --- reading uncommitted state that is later rolled
back --- except the ``read'' is a causal dependency on an external effect.
The harm (invitations already sent) would arise even if those invitations
were themselves reversible.

\paragraph{A7: Conflicting Externalization (source: interaction)}
\emph{Two independent executions externalize non-commuting effects to a
shared resource with no ordering between them, so the outcome depends on
their interleaving.} For example, two coding workers push to the same branch
from the same base:
\begin{verbatim}
effect(push X on main) = confirmed [w1]  w1 pushes from base
effect(push Y on main) = confirmed [w2]  w2 pushes from same base
NOT dep(push Y <- push X)                neither saw the other
NOT commute(push X, push Y)              order changes the result
=> the two effects race on main; outcome depends on interleaving
\end{verbatim}
Unlike contaminated speculation (A6), there is no dependency between the two
executions --- neither observed the other; both acted on the shared resource
independently. The anomaly is neutralized by \emph{commutativity} (if the two
effects commute, order is irrelevant) or by \emph{invertibility} (a versioned
reversible resource lets the loser be rolled back and retried, as in
classical concurrency control); it is irremediable only when the effects
are both non-commuting and irreversible, since atomic externalization across
independent endpoints is impossible without tool-side support.

\paragraph{A8: Phantom Compensation (source: interaction, open-world)}
\emph{An external actor observes and acts on an effect before it is
compensated, so the compensation cannot erase its consequences.} For
example, a supplier reacts to an offer that is later withdrawn:
\begin{verbatim}
effect(send_offer) = confirmed [A]   supplier sees offer
effect(react) = confirmed [X]        supplier acts on it
dep(react <- send_offer)             X depends on the offer
cmp(send_offer) [A]                  offer withdrawn
NOT commute(send_offer, react)       X already acted
=> supplier acted on a now-withdrawn offer; world diverges
\end{verbatim}
Like conflicting externalization, this is a \emph{relational} anomaly: it
turns on commutativity of a pair, not a single operation --- here of the
effect and the external actor's reaction --- which is why it needs a
commutativity \emph{contract} rather than a per-operation hint, and why it is
irremediable when the competitor is the open world (see the
externally-mediated guarantee below).

\section{Effect-Safety Guarantees}
\label{sec:levels}

The anomaly catalog induces a vocabulary of \emph{effect-safety guarantees}:
each guarantee names a set of anomalies it excludes.

\begin{itemize}
\item \textbf{Unknown-safe} (excludes A1, A2, A3): the system never
duplicates, loses a required committed effect, or blindly compensates an
effect whose outcome it cannot confirm.
\item \textbf{Compensation-safe} (excludes A4): an aborted workflow leaves
no irreversible residue, assuming compensations are correct, complete, and
successful.
\item \textbf{Speculation-safe} (excludes A5, A6): no effect that may not
survive is externalized before the workflow is resolved, and no committed
branch depends on an unresolved one.
\item \textbf{Externally-mediated} (excludes A7, A8): interactions with
other executions and outside actors are mediated --- shared-resource writes
are ordered or made to commute, and external reactions to intermediate
effects are controlled --- yielding a serializable external history.
\end{itemize}

The catalog reflects how external the cause of an anomaly is. A1--A6 the
agent creates itself and can in principle prevent through deduplication,
delivery guarantees, outcome resolution, gating, dependency tracking, and
branch isolation. A7 and A8 arise from interaction beyond the execution: a
concurrent execution writing a shared resource (A7), or an outside actor
that observes and reacts to an effect before it is compensated (A8).

Interactive tasks explain why the compensation-safe regime is sometimes the
best achievable one. Consider an agent negotiating on a user's behalf. It
cannot learn the counterparty's reaction without sending an offer, and
sending an offer is already an effect: the action must precede the
observation. Full gating is therefore impossible. The best attainable
guarantee allows the workflow to externalize effects sequentially while
avoiding irreversible residue in discarded branches. When the inputs needed
to resolve the workflow can be observed without effect, for example through
a \texttt{quote} operation, the stronger gated guarantee becomes attainable.

We give two views of these guarantees. The first is a table of
\emph{anomaly $\to$ preventing rule $\to$ assumptions}: a set of ``rules of
the road'', each a local discipline that prevents one anomaly under its
assumptions. The protocols that realize these rules are follow-on work.

\begin{center}
\begin{tabular}{l l p{5.0cm}}
\textbf{Anomaly} & \textbf{Source} & \textbf{Preventing rule (assumptions)} \\
\hline
A1 Duplicate & uncertainty & durable logical-op identity; idempotency key \\
A2 Missing & uncertainty & at-least-once delivery; resolve outcome before commit \\
A3 Orphaned & uncertainty & resolve outcome before compensating \\
A4 Residue & workflow & compensation log + abort barrier (compensations correct, complete) \\
A5 Premature & workflow & stage / gate until resolved \\
A6 Contaminated & workflow & dependency tracking; no commit over unresolved deps \\
A7 Conflicting & interaction & commutativity or invertibility of shared-resource effects; else ordering \\
A8 Phantom & interaction & mediated visibility; observers cannot react to tentative effects \\
\end{tabular}
\end{center}

The second view places existing systems against these guarantees by their
\emph{actual} coverage: no current system realizes a clean cumulative
hierarchy, so we report the set of anomalies each one excludes.

\begin{center}
\renewcommand{\arraystretch}{1.5}
\begin{tabular}{@{}l p{4.6cm} p{3.2cm}@{}}
\textbf{System} & \textbf{Guarantees (anomalies)} & \textbf{Notes} \\
\hline
ACRFence~\cite{zheng2026acrfence} & unknown-safe partial (A1) & proposed; attack validated, not implemented \\
RAC~\cite{perera2026rac} & compensation-safe partial (A4) & no gating; no idempotency, so A1--A3 not addressed \\
Atomix~\cite{mohammadi2026atomix} & A1/A3 partial, compensation-safe (A4), speculation-safe (A5, A6), externally-mediated partial (A7) & idempotent-known-outcome path; no crash-safe exactly-once; A2 gap on multi-effect release \\
Cordon~\cite{chen2026cordon} & A1/A3 partial, compensation-safe (A4), speculation-safe partial (A5) & idempotent tools; staged effects; task-scoped \\
CoAgent~\cite{lyu2026coagent} & externally-mediated partial (A7) & multi-agent serializability via reorder + registered inverse; irreversibles gated \\
\end{tabular}
\renewcommand{\arraystretch}{1}
\end{center}

The coverage entries should be read as the guarantees a system targets under
its stated assumptions, not as proven results. ACRFence targets A1 for the
nondeterministic-replay case: after checkpoint-restore an LLM re-synthesizes
a syntactically different but semantically equivalent irreversible call,
which servers accept as new; ACRFence proposes logging effects and enforcing
replay-or-fork, validating the attack but not yet implementing the
mitigation. RAC targets A4 by compensation rather than gating: effects
externalize \emph{immediately}, and the runtime logs them and runs
compensations on abort in dependency order (dependents before parents).
Because it does not gate, an irreversible effect that has externalized
before an abort can leave residue, and when a compensator is missing RAC
assumes the tool had no side effect; its A4 coverage is therefore partial.
It also has no idempotency mechanism, so the unknown-outcome family (A1--A3)
is not its target.

Atomix and Cordon instead \emph{gate} irreversible
effects until commit, which removes A4 for gated irreversible effects; for
eagerly externalized reversible effects, A4 is excluded under the usual
compensation assumptions. They
avoid A1 and A3 only on paths where unknown outcomes are not blindly
compensated: either the effect is gated and has not externalized, or
re-issue is safe through idempotent, known-outcome tool semantics. This is
not a general resolve-before-compensate protocol --- without idempotency,
status resolution, or durable acknowledgment, A1/A3 fall outside their
guarantee, which is why we mark them partial.

Their remaining gap is also at
the moment of release: several gated irreversibles cannot be externalized
atomically, so a crash mid-release can commit while one required effect is
missing --- an instance of A2, not A4, which is why A2 is absent from their
coverage. Atomix
adds frontier-gated dependency tracking under speculation and contention
(A6), and its frontiers also order concurrent writes to a shared resource,
which partially addresses conflicting externalization (A7) --- but ordering
is not atomicity, so non-commuting irreversible effects can still race;
for Cordon, A7 is outside its task-scoped boundary rather than a guarantee it
targets.

CoAgent addresses A7 directly: it fixes a serialization order at
launch, applies writes speculatively, and mechanically undoes and reorders
misplaced writes through a saga-style inverse each tool registers in advance,
reaching serializability at quiescence. Its two mechanisms are exactly the
gates our catalog predicts --- it reorders when operations commute and inverts
when they do not --- and it hits the same limit: an irreversible effect can be
neither reordered nor inverted, so CoAgent falls back to gating it (its calls
are refused until predecessors commit). It thus realizes externally-mediated
safety only under a registered-inverse (or commutativity) contract, an
instance of the achievability condition for A7 rather than an unconditional
guarantee. No system controls an outside actor's reaction to a released
effect, so phantom compensation (A8) is uncovered.

\section{Research Directions}

Several questions follow from this framework.

\subsection{Achievability}

Given:
\[
G^*(W)
\]
the set of anomalies that workflow $W$ can avoid --- not a point on a scale
but a subset of the catalog, so a \emph{partial} guarantee (some anomalies
excluded, others not) is a first-class value, not a shortfall.

We conjecture:
\[
G^*(W) = F(P, S, C)
\]
where:

\begin{itemize}
\item $P$: per-tool properties (commutativity, invertibility, idempotence,
externalization, determinism);
\item $S$: workflow structure --- ordering, parallelism and speculation,
the dependency graph, and whether read-only (\texttt{quote}) phases precede
effects;
\item $C$: task constraints and invariants (e.g.\ the budget bound).
\end{itemize}

The structure $S$ is not fixed in advance: an agent's execution graph is
determined as it runs. The
agent decides nondeterministically what to call next, so the graph is
\emph{generated dynamically}. We therefore read $S$ not as a single graph
but as the \emph{space} of histories the agent may produce under a given
tool set and policy, and $G^*(W)$ as the set of anomalies avoided across all
of them --- or, equivalently, what a runtime mechanism can enforce as
effects appear, rather than by static analysis of a fixed plan (the
approach taken by systems such as Atomix~\cite{mohammadi2026atomix}). This structural nondeterminism
(\emph{which} operation is invoked) is a distinct axis from the operational
nondeterminism already captured in $P$ (whether a given operation is
deterministic).

Membership in $G^*(W)$ is decided per anomaly: each $A_i$ has a neutralizing
property (Section~\ref{sec:levels}), and $A_i \in G^*(W)$ exactly when $F$
finds that property realizable from $(P, S, C)$ --- idempotence or a
status/ack contract for A1--A3, a correct compensator or a gate for A4,
an invariant checkable without externalizing an effect for A5, dependency
tracking for A6,
commutativity or invertibility on the shared resource for A7, mediated
visibility for A8. A workflow avoids precisely the subset whose conditions
hold, which is why partial coverage is the normal case rather than a
failure to reach a level.

The decisive argument differs by anomaly: unknown-safety follows from per-tool
uncertainty-resolution properties --- idempotency keys, durable
acknowledgments, or status-resolution endpoints --- whereas the gating guarantee against premature
externalization is dominated by the tool \emph{set} and \emph{composition}. For instance, the budget
invariant of \texttt{BookTrip} depends on flight and hotel prices. If
prices are revealed only by booking, verifying the invariant requires
externalizing irreversible effects --- a cycle (to check, book; to book
safely, check) --- and gating is unachievable. If a read-only
\texttt{quote()} exposes prices without effect, the invariant becomes
checkable before any irreversible effect, and gating becomes achievable. This
condition --- that the values an invariant depends on can be observed without
externalization, before irreversible effects must occur --- is necessary for
gating, and in simple staged-effect models, where the only obstacle is the
check-then-act cycle, it is also sufficient.

\subsection{Synthesis: the Inverse Problem}

Achievability asks, given a workflow, which anomalies it can avoid. The inverse
is often what a developer actually wants: given a target set of anomalies to
exclude $A_{target}$
and task constraints $C$, what is the minimal change that makes it
achievable? The change may be a \emph{new tool} (e.g.\ a read-only
\texttt{quote} exposing an invariant's inputs without effect), a
\emph{strengthened tool property} (an idempotency key to reach unknown-safety), or a
\emph{restructuring} (inserting a quote-then-commit phase). Each
corresponds to a move in $P$ or $S$, and the output is effectively a
\emph{specification for a missing tool contract} (see the Tool Contracts
direction below).

Two constraints bound the inverse problem. Minimality must be defined against
\emph{realistic} extensions: a read-only mirror of an effectful operation
is realistic, whereas making an irreversible effect reversible often is
not. And some targets are unachievable by \emph{any} addition --- for
example, interactive tasks where an invariant's input is observable only
through the effect itself (the compensation-safe ceiling). The inverse
problem must therefore distinguish a missing tool from the absence of any
tool that could help.

\subsection{Synthesizing Compensations}

Much of the compensation-based part of the catalog --- A4 and the
recovery of committed effects in commit spheres --- presumes that effects
have compensations. In practice agent
tools rarely ship with them; they arrive as bare API calls. The
achievability of the compensation-based guarantees therefore hinges on whether
a compensator can be \emph{synthesized} for a given operation --- generated
from its signature, description, and observed effect --- together with the
input mapping from the original effect to the compensating call (which
output fields, such as a confirmation reference, the compensation consumes).
This raises two questions the framework can frame. First,
\emph{synthesizability}: for which operations is a correct compensator
derivable, and for which is it provably impossible (a non-refundable
purchase admits none)? This bounds where compensation-safety is reachable at all. Second,
\emph{trust}: a synthesized compensator inherits uncertain semantics
(Section~\ref{sec:unknown}), so its own correctness becomes part of the
guarantee rather than an assumption. We expect synthesizability to be the
practical gate on whether the compensation-based guarantees are inhabited in
real deployments.

\subsection{Commit Spheres}

Some effects cannot be committed independently. In \texttt{BookTrip}, the
charge depends on both bookings (its amount is their sum), and the
confirmation email asserts all three; committing any one without the
others leaves the world inconsistent --- a payment without a booking, or a
confirmation that lies. Such effects must become final together or not at
all. We call a minimal such set a \emph{commit sphere}: the unit of
atomicity for external effects, by analogy with persistence
spheres~\cite{weikum1991principles}, where a minimal set of pages must
reach disk together. (A pure read such as \texttt{search\_flights} has no
effect and belongs to no sphere.)

Commit spheres act at two levels. As \emph{all-or-nothing} fixation they
already matter for compensation-safety and above: a sphere must commit or
abort as a whole, so that crashes and losing branches leave no partial
residue (this is failure atomicity). For external mediation a sphere
additionally requires atomic \emph{visibility}: its effects become visible
to competitors only all at once, so that no one observes an intermediate
state --- which is exactly what would create phantom compensation (A8). Their
size measures the cost of the guarantee: the larger the sphere, the more
coordination, commit latency, and exposure of irreversible effects.

Three features distinguish them from the classical
notion: the execution graph is dynamic, so a sphere is discovered
incrementally as dependent effects appear; its elements are heterogeneous
in reversibility, so an irreversible effect cannot be committed on par with
the rest but must be gated; and the dependency graph may contain cycles
(reflection loops).

\subsection{Completeness and Scope}

The catalog cannot be complete in an absolute sense --- in an open world an
external actor can react in unbounded ways --- so we aim only for completeness
relative to the model, and each family has a structural reason for its
membership. The uncertainty
family is fixed by the outcome of an effect: when it is known, the observed
history can differ from a fault-free run only in multiplicity --- too many
(A1) or too few (A2); when it is unknown, any action taken on it is unsafe
(A3). The workflow family is fixed by how a confirmed effect of a single
execution relates to its resolution point --- through survival under abort
(A4), timing (A5), and dependency (A6), the three internal relations a
single execution admits. The interaction family is fixed by what an effect
meets outside its execution: another execution (A7) or an external actor
(A8), harmful in each case exactly when the two do not commute; the other
combinations reduce to A5/A6, to a stale read, or to information disclosure.

This structure suggests how completeness relative to the model might be
argued: rather than enumerate every possible failure, one would state the
positive lifecycle conditions these families describe and show that a history
satisfying them exhibits none of A1--A8, and conversely. Stating those conditions independently of the anomaly names is
what would keep the equivalence from being a tautology; delimiting the
relation families it rests on (adding authorization, read consistency, or
deadlines would introduce more) is part of the same problem.

We explicitly place the following outside the model, so the catalog is not
mistaken for covering them: semantic planning errors (the agent booked the
wrong city); read-side anomalies (stale or inconsistent reads --- the
subject of read-side isolation work such as S-Bus); security and policy
violations (a secret leaked to a channel); tool-contract misclassification
(a tool declared reversible that is not); the order in which an execution's
own repeated or independent effects are externalized (which a fuller
treatment would add as an uncertainty condition, but which we leave aside
here) --- except where order becomes observable as a shared-resource
interaction between executions, which the catalog does capture, as A7; and
liveness failures (a workflow that never finishes). These may matter, but
they are not external-effect anomalies in the sense studied here.

\subsection{Unknown Operation Semantics}
\label{sec:unknown}

The model so far takes operation properties --- invertibility, idempotence,
commutativity --- as known, as in advanced transaction models. For
LLM-driven agents this is the assumption most often violated: tool contracts
are black boxes, and the agent \emph{infers or guesses} whether an effect is
compensable, idempotent, or commutative, often incorrectly. RAC, for
instance, adopts the simplifying assumption that a tool without a
discoverable compensator has no side effect. This suggests enriching $P$ with an epistemic status for each
property (\emph{known}, \emph{assumed}, \emph{unknown}) and a corresponding
class of anomalies in which an action is taken under a mistaken belief about
semantics --- compensating an operation that has no true inverse, or
replaying one wrongly believed idempotent. Characterizing effect safety under
uncertain semantics, and the safe actions available when a contract is
unknown, is to our knowledge unaddressed, since existing systems take
contracts as given. We see this as the most agent-specific direction the
framework opens.

\subsection{Tool Contracts}

Tool interfaces such as MCP currently expose limited operational
properties. Richer contracts could let a runtime derive achievable
effect-safety guarantees automatically, by declaring the primitives each
preventing rule needs. Several recur across the catalog: a \texttt{quote} or
\texttt{dry\_run} that exposes the inputs of a workflow invariant without
externalizing an effect (enabling the gated guarantee against premature
externalization); a durable idempotency key or stable logical-op identity
(against duplication); a \texttt{hold}/\texttt{reserve} with an explicit
expiry that stages an effect before commitment; a compensation declaration
with its preconditions and an explicit signal when none exists (rather than
silently assuming none is needed); and a status-resolution endpoint that
turns an \emph{unknown} outcome into a definite one (against orphaned
compensation). A contract language expressing these would make the mapping
from tool properties to achievable guarantees mechanical rather than
hand-derived.

The uncertainty family (A1--A3) is a sharp instance of this dependence on
the tool. In the general open-world case it is not resolvable by the runtime
alone: with an unreliable channel the system cannot learn, in bounded time,
whether an effect took place --- the two-generals barrier. Resolvability is
therefore a property of the tool's contract, not of our runtime. An
idempotency key sidesteps the problem (re-issue is safe, so the outcome need
not be known); a status-resolution endpoint or a durable acknowledgment
removes it (the outcome can be learned). With none of these, the family is
irreducible, and a runtime can only choose \emph{which} reaction to risk ---
re-issue, commit optimistically, or compensate --- trading A1 against A2
against A3. Characterizing exactly which tool primitive makes each member
resolvable, and what remains under none, is an open problem and a natural
impossibility result for the protocol layer.

\subsection{Protocols}

This proposal maps the guarantee space; it does not design the protocols
that realize it. The natural next step treats each rule in
Section~\ref{sec:levels} as a protocol to be made precise and proved.
Three questions follow. \emph{Protocol correctness}: prove that a given
protocol prevents a target anomaly class under explicit tool and workflow
assumptions. \emph{Protocol impossibility}: prove that no protocol prevents
a given anomaly without a specific property (e.g.\ that A1 cannot be
prevented without durable logical-op identity, or that A8 cannot be
prevented in an open world without mediated visibility). \emph{Protocol
synthesis}: given a target set of anomalies to exclude, infer the minimal
tool contracts and runtime rules required. This is the constructive
counterpart of the present catalog: a set of rules that prevent classes of
collisions provided all participants comply.

\section{Related Work}

Classical transaction models introduced isolation levels and anomaly
definitions~\cite{berenson1995critique,adya1999}. Advanced transaction
models investigated compensation, sagas, semantic recovery, alternative
execution paths, and dependencies between long-running transactions
~\cite{garcia1987sagas,korth1990formal,chrysanthis1991acta}. Escrow
transactions introduced conditional commutativity and semantic
concurrency control~\cite{oneil1986escrow}. The saga lineage already
isolates the notion of a point of no return --- the pivot after which
effects are no longer compensable --- and prescribes placing irreversible
steps after fallible ones; our catalog can be read as classifying the ways
an effect can be incorrect at that point, extended to concurrent,
speculative, and nondeterministic agent workflows that the single-workflow
saga model does not address.

Recent systems build runtime mechanisms for agent effects; we read them as
\emph{partial protocol implementations} for individual rules in
Section~\ref{sec:levels}, and discuss their mechanisms there. ACRFence
~\cite{zheng2026acrfence} proposes a mitigation against duplicated effects
under nondeterministic replay (validated as an attack, not yet implemented);
RAC~\cite{perera2026rac} is a compensation protocol against abort
residue; Atomix~\cite{mohammadi2026atomix} and Cordon~\cite{chen2026cordon}
implement staging/gating under a mediated boundary; and
CoAgent~\cite{lyu2026coagent} builds multi-agent concurrency control on a
shared resource, reordering or inverting conflicting writes to reach
serializability --- the closest system to our conflicting-externalization
anomaly (A7). SagaLLM~\cite{chang2025sagallm} applies Saga-style validation to
multi-agent planning, and Stonebraker et
al.~\cite{stonebraker2025consistency} call for isolation-level reasoning in
data-oriented workflow systems.

Cordon introduces semantic transactions --- task-level boundaries tracking
tool intents, result lineage, reversible state, staged effects, and
authority --- which closely mirror our effect histories, dependency edges,
and stage/gate distinction; but it targets security and least-privilege
containment rather than a general anomaly-based account of external-effect
guarantees under concurrency, recovery, and speculation. On a complementary
axis,
S-Bus~\cite{khan2026sbus} formalizes \emph{read}-side isolation: its
Observable-Read Isolation governs what an agent observes of shared mutable
state, where we govern what it irreversibly externalizes --- the read-side
and effect-side of the same problem. On another complementary axis,
Zhai et al.~\cite{zhai2026revisable} study \emph{intent revision} during
streaming execution, classifying actions by reversibility and showing that
conflicting irreversible actions make full specification satisfaction
impossible --- an algorithm-independent limit that parallels our own
achievability bounds; their earliest-conflict rollback, however, is defined
over a linear action sequence, whereas our dependency relation is a partial
order, under which ``earliest'' is not uniquely defined.

\bibliographystyle{plain}
\bibliography{bibliography/flame-stream}

\end{document}